\section{A variational bound retaining both real zeros and the known prime}
\label{sec:s5-variational-endpoint}

This refines the residual interval in MA-RB-002 without a
zero-location hypothesis.  It uses the exact signature and the
specified Frobenius at five, rather than replacing all their
contributions by a constant.  Throughout, \(L/\mathbb Q\) is
the actual Galois closure in Calegari's family and
\(\Delta=\Delta_L\) is its absolute discriminant.
Thus \([L:\mathbb Q]=120\), its signature is \((0,60)\),
and five is unramified with Frobenius of cycle type \(2B\).

\paragraph{MA-RB-007 (the exact one-variable interval bound).}
For \(\sigma>1\), define
\[
 \begin{split}
 H_\Delta(\sigma)={}&\frac1\sigma+\frac1{\sigma-1}
       +\frac12\log\Delta
       +60\bigl(\psi(\sigma)-\log(2\pi)\bigr)\\
       &-\frac{120\log5}{5^{2\sigma}-1},
       \qquad \psi=\Gamma'/\Gamma .
 \end{split}
 \tag{MA-RB.71}
\]
For every actual field in this family, \(H_\Delta(\sigma)>0\)
on the entire indicated half-line.  For an integer \(k\ge1\), put
\[
 \begin{aligned}
 F_{\Delta,k}(\sigma)
    &=(\sigma-\tfrac12)^2
          -\frac{9k(2\sigma-1)}{H_\Delta(\sigma)},\\
 J_{\Delta,k}&=\inf_{\sigma>1}F_{\Delta,k}(\sigma).
 \end{aligned}
 \tag{MA-RB.72}
\]
This is a finite real number, strictly less than \(1/4\),
and its infimum is attained at a point of \((1,\infty)\).

Suppose the remaining comparison-failure branch of
MA-RB-002 occurs.  Let \(\gamma\in(1/2,1)\) be its
rightmost real abscissa and \(k\ge1\) the order of the
comparison pole there.  Then
\[
 \begin{gathered}
 9k\left(\frac1{\sigma-\gamma}
               +\frac1{\sigma-1+\gamma}\right)
                  \le H_\Delta(\sigma)\qquad(\sigma>1),\\
             (\gamma-\tfrac12)^2\le J_{\Delta,k}.
 \end{gathered}
 \tag{MA-RB.73}
\]
Consequently \(J_{\Delta,k}\le0\) excludes that branch;
otherwise its exact bound from these retained terms is
\[
             \frac12<\gamma
                  \le\frac12+\sqrt{J_{\Delta,k}}<1.
 \tag{MA-RB.74}
\]
Since \(J_{\Delta,k}\le J_{\Delta,1}\), the criterion
\(J_{\Delta,1}\le0\) excludes the comparison-failure
branch without knowing its pole order.  If
\(J_{\Delta,1}>0\), the residual shared-zero criterion
in MA-RB-002 can be confined to
\[
 \begin{gathered}
 B(\gamma)=0,\qquad L(\gamma,V\eta)L(\gamma,X)=0,\\
 \frac12<\gamma\le\frac12+\sqrt{J_{\Delta,1}}.
 \end{gathered}
 \tag{MA-RB.75}
\]
These statements do not assert that either the numerical
criterion or the remaining nonvanishing condition holds
for every discriminant in the family.

\begin{proof}
\emph{The known finite-prime term and positivity of the denominator.}
The order of a \(2B\) element is two.  Since five is
unramified and \(L/\mathbb Q\) is Galois of degree 120,
every prime above five has residue degree two and there
are exactly sixty such primes.  Its complete contribution
to the positive Euler logarithmic derivative is therefore
\[
 E_5(\sigma)
    =60\sum_{j\ge1}\frac{\log25}{25^{j\sigma}}
    =\frac{120\log5}{5^{2\sigma}-1}.
 \tag{MA-RB.76}
\]
In particular this retains all powers of five, not only
the first local coefficient.

The completed Dedekind function and the constant in its
Hadamard logarithmic derivative were fixed and proved in
MA-RB.8--.12.  They give the exact, absolutely convergent
identity
\[
 \begin{split}
 H_\Delta(\sigma)
  ={}&\sum_\rho\operatorname{Re}\frac1{\sigma-\rho}\\
     &+\left(-\frac{\zeta_L'}{\zeta_L}(\sigma)
                                      -E_5(\sigma)\right).
 \end{split}
 \tag{MA-RB.77}
\]
The sum includes multiplicities of all nontrivial completed
zeros.  The unconditional containing strip
\(0\le\operatorname{Re}\rho\le1\) makes every summand
nonnegative for \(\sigma>1\).  The parenthesized term
is the sum of the strictly positive Euler terms at all
prime ideals not above five.  It is strictly positive,
for example because there is a prime ideal above two.
This proves \(H_\Delta>0\) for an actual field, independently
of the proposed comparison failure or the existence of
any specified zero.  Only the usual containing strip is
used; no critical-line location is used.

\emph{The two retained real zeros.}
MA-RB-002 proves that \(\zeta_L\) has a zero of
multiplicity at least \(9k\) at \(\gamma\).
The completed functional equation
\(\xi_L(s)=\xi_L(1-s)\) gives the same multiplicity
at \(1-\gamma\).  These points are distinct, lie in
\((0,1)\), and none of the completion factors has a
zero or pole there.  Thus both occurrences belong to
the nontrivial completed zero divisor in (MA-RB.77).
Retaining them and dropping only nonnegative other
terms gives the first assertion of (MA-RB.73).

Write \(a=\gamma-1/2\) and \(x=\sigma-1/2\).
Then \(0<a<1/2<x\), and their contribution is
\[
 9k\left(\frac1{x-a}+\frac1{x+a}\right)
          =\frac{18kx}{x^2-a^2}.
 \tag{MA-RB.78}
\]
All denominators are positive.  Multiplying the first
inequality of (MA-RB.73) by them gives
\(a^2\le x^2-18kx/H_\Delta(\sigma)\).
Taking the infimum over \(\sigma>1\) gives exactly
the second assertion.  The conclusions for nonpositive
and positive \(J\) follow, and monotonicity in \(k\)
follows directly from the positive denominator in
(MA-RB.72).  MA-RB-002 already identifies the two
shared vanishing factors at this actual \(\gamma\);
only their permitted interval is changed in (MA-RB.75).

\emph{Existence and range of the variational minimum.}
All terms in \(H_\Delta(1+t)-1/t\) are smooth at
\(t=0\).  Hence
\[
 \begin{aligned}
 H_\Delta(1+t)&=t^{-1}+O_\Delta(1),\\
 F_{\Delta,k}(1+t)
       &=\frac14+(1-9k)t+O_{\Delta,k}(t^2)
                         \qquad(t\downarrow0).
 \end{aligned}
 \tag{MA-RB.79}
\]
The coefficient of \(t\) is negative.  Thus the
continuous extension of \(F\) at \(\sigma=1\) has
value \(1/4\), and \(F\) is strictly below that
value at some interior point.
For completeness, take a gamma random variable \(T\)
of shape \(\sigma\) and rate one.  As in MA-RB-001,
\(\psi(\sigma)=\mathbb E\log T\) and
\(\mathbb ET=\sigma\).  When \(\sigma>1\), the gamma
recurrence also gives \(\mathbb E(T^{-1})=(\sigma-1)^{-1}\).
Jensen's inequality applied to \(T\) and \(T^{-1}\)
therefore gives
\(\log(\sigma-1)\le\psi(\sigma)\le\log\sigma\).
In particular \(\psi(\sigma)=\log\sigma+O(1/\sigma)\),
which implies
\[
 H_\Delta(\sigma)=60\log\sigma+O_\Delta(1),\qquad
 F_{\Delta,k}(\sigma)\longrightarrow+\infty
                       \quad(\sigma\longrightarrow\infty).
 \tag{MA-RB.80}
\]
On every compact subinterval of \((1,\infty)\),
\(H\) is positive and continuous, so \(F\) is finite
and continuous.  These facts confine the infimum to
a compact subinterval strictly inside \((1,\infty)\).
It is attained there, is finite, and is strictly below
\(1/4\).  This proves all assertions about the minimum.
\end{proof}

\paragraph{Strict improvement over the previous endpoint.}
The variational restriction is strictly stronger than the
one-sided bound of MA-RB.21 whenever the hypothetical
branch is feasible.  Here is a direct comparison that
does not require numerically finding the minimizer.
Put
\[
 \begin{gathered}
 \mathcal H_L=1+\tfrac12\log\Delta,\\
 t_0=\frac{3\sqrt{k}-1}{\mathcal H_L},\qquad
 \delta_0=\frac{(3\sqrt{k}-1)^2}{\mathcal H_L}.
 \end{gathered}
 \tag{MA-RB.81}
\]
Under that branch, MA-RB-002 proves \(0<t_0<1\)
and \(0<\delta_0\le1-\gamma<1/2\).
The exact old optimization gives
\[
              \frac{9k}{t_0+\delta_0}
                            =\frac1{t_0}+\mathcal H_L.
 \tag{MA-RB.82}
\]
For \(1<\sigma\le2\), the gamma-integral Jensen
bound proved in MA-RB-001 gives
\(\psi(\sigma)\le\log\sigma<\log(2\pi)\).
Also \(1/\sigma<1\), and \(E_5(\sigma)>0\).
Consequently
\[
 \begin{split}
 H_\Delta(1+t_0)&<\frac1{t_0}+\mathcal H_L\\
       &<9k\left(\frac1{t_0+\delta_0}
                         +\frac1{1+t_0-\delta_0}\right).
 \end{split}
 \tag{MA-RB.83}
\]
The last expression is (MA-RB.78) evaluated at
\(x=1/2+t_0\) and \(a_0=1/2-\delta_0>0\).
The same positive-denominator algebra now gives
\[
 J_{\Delta,k}\le F_{\Delta,k}(1+t_0)
                         <(\tfrac12-\delta_0)^2.
 \tag{MA-RB.84}
\]
If the branch survives, \(J>0\), so its new upper
endpoint is strictly smaller than \(1-\delta_0\).
If \(J\le0\), it has instead been excluded.
For a pole order not yet known, the choice \(k=1\)
therefore gives a strict refinement of the universal
endpoint \(1-4/\mathcal H_L\) whenever that branch
had remained feasible.

\paragraph{Source and method scope.}
The prime-five hypothesis is Calegari's Theorem~1.1,
assumption~(2), in the supplied \texttt{S5ArtinWeb.tex}
lines 123--136; assumption~(1) supplies the totally
imaginary signature.  The class \((12)(34)\) has
order two; its label \(2B\) follows the checked table.
The completed Hadamard identity, including its additive
constant and all sixty complex gamma factors, is the
already proved MA-RB.8--.12.  The multiplicity \(9k\)
and the shared-zero reduction are MA-RB-002.
The present optimization is a deduction from those
exact inputs, not a numerical test of a field or an
additional automorphy or zero-location theorem.
It does not use \(\mathsf A\), \(M\), or
\((\mathcal R_B)\).
